\documentclass[11pt]{article}
\usepackage[margin=1.1in]{geometry}
\usepackage{amsmath, amssymb, amsthm}
\usepackage{mathtools}
\usepackage{booktabs}
\usepackage{xcolor}
\usepackage{hyperref}
\usepackage{microtype}
\usepackage{parskip}

\hypersetup{colorlinks=true, linkcolor=blue!60!black, urlcolor=blue!60!black}

\definecolor{abilityblue}{HTML}{2C5F8A}
\definecolor{willgreen}{HTML}{3A7A4E}

\title{\textbf{``Can you?'' as Indirect Request}\\[4pt]
\large A Rational Speech Acts Model: Ability-Only and Willingness Extension}
\author{can\_you project}
\date{}

\begin{document}
\maketitle
\tableofcontents
\bigskip

% ============================================================
\section{Overview}
% ============================================================

The sentence \emph{``Can you pass the salt?''} is grammatically a yes/no question
about ability, yet speakers routinely use it as a request. This document presents
two versions of a \textbf{Rational Speech Acts (RSA)} model~\cite{frank2012predicting,goodman2016pragmatic}
that explain---without ad hoc context-specific stipulations, and derived from
the literal semantics of ``can you?'' and standard RSA assumptions---when a
listener should interpret \emph{``Can you?''} as a request versus a genuine
information-seeking question, and how that interpretation drives whether the
listener acts.

\begin{itemize}
  \item \textbf{Ability-only model} (\texttt{whyask\_ability.ipynb},
        \texttt{test\_minimal.py}): compliance depends solely on ability
        ($\mu > \theta$). One context parameter: $f_{\text{ctx}}$ (feasibility).
  \item \textbf{Willingness model} (\texttt{whyask\_with\_willingness.ipynb}):
        compliance requires ability \emph{and} willingness
        (($\mu > \theta$) \emph{and} ($v > \phi$)). Two context parameters:
        $f_{\text{ctx}}$ (ability) and $w_{\text{ctx}}$ (willingness).
\end{itemize}

The core insight shared by both: \textbf{context determines the informational
value of asking.} When ability is already common ground (everyone knows you can
pass the salt), the question carries almost no information, so a pragmatic
listener infers the speaker must have had a different motive.

Both models use two levels of speaker/addressee recursion (S1/A1 and S2/A2),
implemented in JAX with the \texttt{memo} probabilistic DSL.

% ============================================================
\section{Shared Setup}
% ============================================================

\subsection{Utterances, goals, responses}

\[
  \mathcal{U} = \{\,\textsc{can},\ \textsc{imp},\ \textsc{null}\,\}
  \qquad
  \mathcal{G} = \{\,\textsc{info},\ \textsc{action}\,\}
  \qquad
  \mathcal{R} = \{\,\textsc{act},\ \textsc{pass}\,\}
\]

\textsc{can} = ``Can you?''; \textsc{imp} = direct imperative (``Pass the salt'');
\textsc{null} = silence. \textsc{info} goal = learn whether addressee can/will act;
\textsc{action} goal = get them to do it.

\subsection{Latent ability variables}

Two continuous variables govern whether the addressee can comply:
\begin{itemize}
  \item $\mu \in [0,1]$: the addressee's latent ability level.
  \item $\theta \in [0,1]$: the implicit threshold for counting as ``able.'' The addressee is able iff $\mu > \theta$.
\end{itemize}

Both are marginalized out in the final predictions (see Sections~\ref{sec:ability} and~\ref{sec:willingness}).
The willingness extension adds two analogous variables $v$ (willingness level) and $\phi$ (willingness threshold).

\subsection{Information gain}

The literal semantics of \textsc{can} is: \emph{``Is $\mu > \theta$?''}---a
binary question about ability. Let $Y = \mathbf{1}[\mu > \theta]$ be the binary
answer (1 = able, 0 = not able). The utility of asking \textsc{can} for an
\textsc{info}-goal speaker is the mutual information between $\mu$ and $Y$:
\[
  \mathrm{IG}(f_{\text{ctx}}, \theta)
  = H\!\left[\pi(\mu)\right]
  - P(Y{=}1)\, H\!\left[\pi(\mu \mid Y{=}1)\right]
  - P(Y{=}0)\, H\!\left[\pi(\mu \mid Y{=}0)\right]
\]
where $\pi(\mu) \equiv \pi(\mu \mid f_{\text{ctx}})$ is a Beta prior over ability
(defined in Section~\ref{sec:ability}), $H[\cdot]$ is Shannon entropy in nats,
and the posteriors are:
\[
  \pi(\mu \mid Y{=}1) \;\propto\; \pi(\mu)\cdot\mathbf{1}[\mu > \theta]
  \qquad
  \pi(\mu \mid Y{=}0) \;\propto\; \pi(\mu)\cdot\mathbf{1}[\mu \leq \theta]
\]
This function is identical in both model versions---\textsc{can} is always
framed as an ability question grammatically.

\textbf{Intuition:} When $f_{\text{ctx}} \approx 1$, the prior over $\mu$ is
concentrated near 1; the answer is unsurprising---IG $\approx 0$. When
$f_{\text{ctx}}$ is low, the answer is genuinely uncertain---IG is high.

\subsection{Response semantics: acting entails yes}

\textbf{Acting answers the ability question.} If you pass the salt, you've implicitly
confirmed you can. This motivates a single response utility function (shared by both models):
\[
  \mathrm{EU}_{\text{respond}}(r, g, \ldots) =
  \begin{cases}
    \text{comply} & \text{if } r = \textsc{act} \\
    \mathbf{1}[g = \textsc{info}] & \text{if } r = \textsc{pass}
  \end{cases}
\]
where ``comply'' is defined differently across model versions (see below).
Passing is good only for an info-goal speaker (truthful verbal answer); acting
requires the addressee to actually be able (and willing). Zero new parameters.

\subsection{Speaker social cost}

The imperative incurs a flat face-threat cost $\delta = 0.7$ regardless of context.
\textsc{can} and \textsc{null} are costless (or near-costless in the notebook
version: $c_{\text{can}} = 0.1$).

\subsection{Softmax choice}

Both speakers and addressees choose via softmax with rationality $\alpha = 10$:
\[
  P(x \mid \ldots) \;\propto\; \exp\!\left(\alpha \cdot \mathrm{EU}(x, \ldots)\right)
\]

% ============================================================
\section{Ability-Only Model}
% ============================================================
\label{sec:ability}

\subsection{Setup additions}

The latent variables $\mu$ and $\theta$ are as defined in Section~2.2.

\paragraph{Feasibility prior.}
\[
  \mu \sim \mathrm{Beta}\!\left(f_{\text{ctx}}\cdot k,\; (1 - f_{\text{ctx}})\cdot k\right), \quad k = 10
\]
Single context parameter $f_{\text{ctx}} \in [0,1]$.

\paragraph{Goal prior.} Action goals presuppose feasibility; given feasibility,
max-entropy:
\[
  P(g = \textsc{action} \mid f_{\text{ctx}}) = \tfrac{1}{2}\,f_{\text{ctx}}
  \qquad
  P(g = \textsc{info} \mid f_{\text{ctx}}) = 1 - \tfrac{1}{2}\,f_{\text{ctx}}
\]

\paragraph{Compliance.}
\[
  \text{comply} = \mathbf{1}[\mu > \theta]
\]

\paragraph{Test contexts.}

\begin{center}
\renewcommand{\arraystretch}{1.2}
\begin{tabular}{lc}
\toprule
\textbf{Context} & $f_{\text{ctx}}$ \\
\midrule
Benchpress 150 lbs & 0.10 \\
Drive me home      & 0.70 \\
Call me later      & 0.90 \\
Pass the salt      & 0.99 \\
\bottomrule
\end{tabular}
\end{center}

\subsection{Speaker utility (S1)}

\[
  \mathrm{EU}_{S1}(u, g, \mu, \theta, f_{\text{ctx}}) =
  \underbrace{\mathbf{1}[g = \textsc{action}]\cdot \mathbf{1}[u = \textsc{imp}]\cdot\mathbf{1}[\mu > \theta]}_{\text{action utility}}
  + \underbrace{\mathbf{1}[g = \textsc{info}]\cdot \mathbf{1}[u = \textsc{can}]\cdot\mathrm{IG}(f_{\text{ctx}}, \theta)}_{\text{info utility}}
  - \underbrace{\delta\cdot\mathbf{1}[u = \textsc{imp}]}_{\text{social cost}}
\]

S1 knows goal $g$, marginalizes over $(\mu, \theta)$:
\[
  P_{S1}(u \mid g, f_{\text{ctx}}) \;\propto\;
  \exp\!\Bigl(\alpha\cdot \mathbb{E}_{\mu \sim \pi(\cdot|f_{\text{ctx}}),\, \theta \sim \mathrm{Unif}}\!\bigl[\mathrm{EU}_{S1}(u, g, \mu, \theta, f_{\text{ctx}})\bigr]\Bigr)
\]

\subsection{Level-1 addressee (A1)}

A1 observes $u = \textsc{can}$ and updates on S1's distribution:
\[
  P_{A1}(g \mid u = \textsc{can}) \;\propto\; P(g \mid f_{\text{ctx}})\cdot P_{S1}(\textsc{can} \mid g, f_{\text{ctx}})
\]

Let $p_a \coloneqq P_{A1}(g{=}\textsc{action} \mid u{=}\textsc{can})$ denote the posterior probability that the speaker had an action goal. A1's response given known $(\mu, \theta)$:
\[
  P_{A1}(\textsc{act} \mid \mu, \theta) =
  \frac{\exp(\alpha\cdot\mathbf{1}[\mu > \theta])}
       {\exp(\alpha\cdot\mathbf{1}[\mu > \theta]) + \exp(\alpha\cdot(1 - p_a))}
\]

Aggregate compliance (marginalized):
\[
  P_{A1}(\textsc{act} \mid f_{\text{ctx}}) =
  \mathbb{E}_{\mu \sim \pi(\cdot|f_{\text{ctx}})}\,
  \mathbb{E}_{\theta \sim \mathrm{Unif}}\!\left[P_{A1}(\textsc{act} \mid \mu, \theta)\right]
\]

\subsection{Level-2 speaker (S2)}

S2 anticipates A1's compliance when using \textsc{can} with an action goal:
\[
  u_{\text{act}}^{S2}(u, f_{\text{ctx}}) =
  \begin{cases}
    P_{A1}(\textsc{act} \mid f_{\text{ctx}}) & u = \textsc{can} \\
    \mathbf{1}[\mu > \theta] & u = \textsc{imp} \\
    0 & u = \textsc{null}
  \end{cases}
\]

Everything else (info utility, $\delta$) is unchanged. S2 can use \textsc{can}
to get compliance while avoiding the face cost of \textsc{imp}---but only in
high-$f_{\text{ctx}}$ contexts where A1 is likely to comply.

\subsection{Level-2 addressee (A2)}

A2 has the same structure as A1 but conditions on S2:
\[
  P_{A2}(g \mid u = \textsc{can}) \;\propto\; P(g \mid f_{\text{ctx}})\cdot P_{S2}(\textsc{can} \mid g, f_{\text{ctx}})
\]
A2's response is computed identically to A1's using the updated posterior.

% ============================================================
\section{Willingness Extension}
% ============================================================
\label{sec:willingness}

\subsection{Motivation}

The ability-only model collapses two distinct reasons not to act: \emph{can't}
vs.\ \emph{won't}. Consider:
\begin{itemize}
  \item \emph{``Can you benchpress 150?''} --- low ability, high willingness
  \item \emph{``Can you give me \$50?''} --- high ability, low willingness
\end{itemize}
Both have low compliance, but for different reasons. The willingness extension
separates these dimensions, enabling a richer 2D context space.

\subsection{New latent variables}

\begin{center}
\renewcommand{\arraystretch}{1.2}
\begin{tabular}{lll}
\toprule
\textbf{Variable} & \textbf{Range} & \textbf{Meaning} \\
\midrule
$\mu$ & $[0,1]$ & addressee's ability level \\
$\theta$ & $[0,1]$ & ability threshold (``able'' iff $\mu > \theta$) \\
$v$ & $[0,1]$ & addressee's willingness level \\
$\phi$ & $[0,1]$ & willingness threshold (``willing'' iff $v > \phi$) \\
\bottomrule
\end{tabular}
\end{center}

The ability prior is as before: $\mu \sim \mathrm{Beta}(f_{\text{ctx}}\cdot k,\, (1-f_{\text{ctx}})\cdot k)$.
Willingness has the same structure: $v \sim \mathrm{Beta}(w_{\text{ctx}}\cdot k,\, (1-w_{\text{ctx}})\cdot k)$.

\subsection{Changes to compliance and goal prior}

\paragraph{Compliance.} Both ability \emph{and} willingness must exceed threshold:
\[
  \text{comply} = \mathbf{1}[\mu > \theta]\cdot\mathbf{1}[v > \phi]
\]

\paragraph{Goal prior.} Action goals presuppose \emph{both} ability and willingness:
\[
  P(g = \textsc{action} \mid f_{\text{ctx}}, w_{\text{ctx}}) = \tfrac{1}{2}\,f_{\text{ctx}}\cdot w_{\text{ctx}}
\]
This means the action goal is only likely when the addressee is both plausibly
able \emph{and} plausibly willing.

\subsection{Speaker utility (S1, willingness version)}

\[
  \mathrm{EU}_{S1}(u, g, \mu, \theta, v, \phi, f_{\text{ctx}}) =
  \mathbf{1}[g{=}\textsc{action}]\cdot\mathbf{1}[u{=}\textsc{imp}]\cdot\mathbf{1}[\mu{>}\theta]\cdot\mathbf{1}[v{>}\phi]
  + \mathbf{1}[g{=}\textsc{info}]\cdot\mathbf{1}[u{=}\textsc{can}]\cdot\mathrm{IG}(f_{\text{ctx}}, \theta)
  - \delta\cdot\mathbf{1}[u{=}\textsc{imp}]
\]

S1 now marginalizes over four latent variables $(\mu, \theta, v, \phi)$:
\[
  P_{S1}(u \mid g, f_{\text{ctx}}, w_{\text{ctx}}) \;\propto\;
  \exp\!\Bigl(\alpha\cdot
    \mathbb{E}_{\mu,\theta,v,\phi}\!\bigl[\mathrm{EU}_{S1}(u, g, \mu, \theta, v, \phi, f_{\text{ctx}})\bigr]
  \Bigr)
\]

\paragraph{Note on information gain.} IG still only depends on $f_{\text{ctx}}$ and
$\theta$ (not on $v$ or $\phi$), because \textsc{can} is grammatically an
ability question. Adding willingness changes the action payoff but not the
literal semantics of the question.

\subsection{A1, S2, A2 (willingness version)}

The structure is identical to the ability-only model; only the definitions of
comply and goal prior change. Goal inference via Bayes:
\[
  P_{A1}(g \mid u = \textsc{can}) \;\propto\; P(g \mid f_{\text{ctx}}, w_{\text{ctx}})\cdot P_{S1}(\textsc{can} \mid g, f_{\text{ctx}}, w_{\text{ctx}})
\]

Response utility uses the willingness-aware comply. As before, let
$p_a \coloneqq P_{A1}(g{=}\textsc{action} \mid u{=}\textsc{can})$ (now computed
from the willingness-aware goal posterior above). Then:
\[
  P_{A1}(\textsc{act} \mid \mu, \theta, v, \phi) =
  \frac{\exp\!\bigl(\alpha\cdot\mathbf{1}[\mu{>}\theta]\cdot\mathbf{1}[v{>}\phi]\bigr)}
       {\exp\!\bigl(\alpha\cdot\mathbf{1}[\mu{>}\theta]\cdot\mathbf{1}[v{>}\phi]\bigr) + \exp(\alpha\cdot(1-p_a))}
\]

Aggregate compliance now marginalizes over four dimensions:
\[
  P_{A1}(\textsc{act} \mid f_{\text{ctx}}, w_{\text{ctx}}) =
  \mathbb{E}_{\mu, \theta, v, \phi}\!\left[P_{A1}(\textsc{act} \mid \mu, \theta, v, \phi)\right]
\]

\subsection{Test contexts}

Sorted by $f_{\text{ctx}} \times w_{\text{ctx}}$ (the product governs goal prior):

\begin{center}
\renewcommand{\arraystretch}{1.2}
\begin{tabular}{lccc}
\toprule
\textbf{Context} & $f_{\text{ctx}}$ & $w_{\text{ctx}}$ & $f \times w$ \\
\midrule
Run a marathon   & 0.10 & 0.01 & 0.001 \\
Benchpress 150   & 0.10 & 0.90 & 0.090 \\
Give me \$50     & 0.80 & 0.20 & 0.160 \\
Drive me home    & 0.70 & 0.80 & 0.560 \\
Call me later    & 0.90 & 0.80 & 0.720 \\
Pass the salt    & 0.99 & 0.99 & 0.980 \\
\bottomrule
\end{tabular}
\end{center}

These contexts reveal dissociations invisible to the ability-only model:
\emph{``Run a marathon''} and \emph{``Benchpress 150''} have the same $f_{\text{ctx}} = 0.10$
but differ dramatically in $w_{\text{ctx}}$, allowing the model to predict
different levels of action inference and compliance even when ability is equally
low.

% ============================================================
\section{Model Comparison}
% ============================================================
\label{sec:comparison}

\subsection{What changes}

\begin{center}
\renewcommand{\arraystretch}{1.3}
\begin{tabular}{lll}
\toprule
\textbf{Component} & \textbf{Ability-only} & \textbf{Willingness} \\
\midrule
Context params   & $f_{\text{ctx}}$ & $f_{\text{ctx}},\, w_{\text{ctx}}$ \\
Latent variables & $\mu, \theta$ & $\mu, \theta, v, \phi$ \\
Compliance       & $\mathbf{1}[\mu > \theta]$ & $\mathbf{1}[\mu > \theta]\cdot\mathbf{1}[v > \phi]$ \\
Goal prior       & $\frac{1}{2}f_{\text{ctx}}$ & $\frac{1}{2}f_{\text{ctx}} w_{\text{ctx}}$ \\
IG               & $\mathrm{IG}(f_{\text{ctx}}, \theta)$ & same \\
S1 utility       & marginalize over $(\mu, \theta)$ & marginalize over $(\mu, \theta, v, \phi)$ \\
A1 response      & marginalize over $(\mu, \theta)$ & marginalize over $(\mu, \theta, v, \phi)$ \\
Number of params & 3 ($\alpha, \delta, k$) & 3 ($\alpha, \delta, k$; same values) \\
Test contexts    & 4, ordered by $f$ & 6, ordered by $f \times w$ \\
\bottomrule
\end{tabular}
\end{center}

No new free parameters are added in the willingness extension---willingness enters
only through the context specification $(f_{\text{ctx}}, w_{\text{ctx}})$, which is
set externally per context.

\subsection{What stays the same}

\begin{itemize}
  \item Information gain: still $\mathrm{IG}(f_{\text{ctx}}, \theta)$---the question
        is grammatically about ability regardless of model version.
  \item RSA structure: S1 $\to$ A1 $\to$ S2 $\to$ A2 recursion is identical.
  \item Social cost: $\delta = 0.7$ flat for \textsc{imp}.
  \item Softmax rationality: $\alpha = 10$.
  \item ``Acting entails yes'' response semantics.
\end{itemize}

\subsection{Key predictions and dissociations}

\paragraph{Ability-only model.} Both \emph{``Benchpress 150''} and
\emph{``Run a marathon''} have $f_{\text{ctx}} = 0.10$, so the model assigns
them identical predictions. It cannot distinguish between ``I won't'' and
``I can't.''

\paragraph{Willingness model.} These contexts now diverge: \emph{``Benchpress 150''}
($f{=}0.10, w{=}0.90$) has much higher $P(g{=}\textsc{action})$ than
\emph{``Run a marathon''} ($f{=}0.10, w{=}0.01$), because asking about something
the addressee is willing but unable to do still makes more pragmatic sense as a
request than asking about something they're both unable and unwilling to do.

Similarly, \emph{``Give me \$50''} ($f{=}0.80, w{=}0.20$) has moderate ability but
low willingness: asking has informational value (the addressee is able but might
not want to), so action inference is lower than \emph{``Drive me home''}
($f{=}0.70, w{=}0.80$) despite similar ability.

\paragraph{A2 amplification.} In both models, A2 infers action intent more strongly than A1 in
high-$f_{\text{ctx}}$ (or high-$f_{\text{ctx}} w_{\text{ctx}}$) contexts
(all quantities conditioned on observing $u = \textsc{can}$).
From the willingness notebook:

\begin{center}
\renewcommand{\arraystretch}{1.2}
\begin{tabular}{lccccc}
\toprule
\textbf{Context} & $f$ & $w$ & $P(g{=}\textsc{action} \mid \textsc{can}, A1)$ & $P(g{=}\textsc{action} \mid \textsc{can}, A2)$ & $P(\textsc{act} \mid \textsc{can}, A2)$ \\
\midrule
Run a marathon   & 0.10 & 0.01 & 0.000 & 0.000 & 0.001 \\
Benchpress 150   & 0.10 & 0.90 & 0.031 & 0.039 & 0.055 \\
Give me \$50     & 0.80 & 0.20 & 0.048 & 0.070 & 0.111 \\
Drive me home    & 0.70 & 0.80 & 0.114 & 0.306 & 0.524 \\
Call me later    & 0.90 & 0.80 & 0.131 & 0.441 & 0.696 \\
Pass the salt    & 0.99 & 0.99 & 0.149 & 0.613 & 0.862 \\
\bottomrule
\end{tabular}
\end{center}

The A2 amplification is largest in the high-$f \times w$ regime---where S2 has
the most to gain from using \textsc{can} strategically.

% ============================================================
\section{The \texttt{memo} DSL}
% ============================================================

\subsection{Core syntax}

\begin{center}
\renewcommand{\arraystretch}{1.4}
\begin{tabular}{ll}
\toprule
\textbf{Syntax} & \textbf{Meaning} \\
\midrule
\texttt{@memo} & marks function as probabilistic computation \\
\texttt{[\,\_x: X,\,\ldots\,]} & indexed return dims (marginalized grid) \\
\texttt{agent: knows(x)} & conditions on agent knowing $x$ \\
\texttt{agent: given(x in X, wpp=...)} & prior over $x$ (wpp = ``with prob.\ proportional'') \\
\texttt{agent: chooses(u in U, wpp=...)} & softmax choice \\
\texttt{agent: observes\_that[cond]} & Bayesian update on observed event \\
\texttt{agent: thinks[\,inner\,]} & agent models inner agent \\
\texttt{Pr[cond]} & return marginal probability \\
\texttt{EU[payoff]} & expected utility, marginalizing over uncertain vars \\
\bottomrule
\end{tabular}
\end{center}

\subsection{S1 (ability-only vs.\ willingness)}

Ability-only (\texttt{whyask\_ability.ipynb}):
\begin{verbatim}
@memo
def S1[_u: U, _g: goals](f_ctx, alpha):
    speaker: knows(_g)
    speaker: given(mu in mu_vals, wpp=f_prior(mu, f_ctx))
    speaker: given(t in theta_vals, wpp=1)
    speaker: wants(payoff = eu_s1(u, _g, mu, t, f_ctx))
    speaker: chooses(u in U, wpp=exp(alpha * EU[payoff]))
    return Pr[speaker.u == _u]
\end{verbatim}

Willingness extension (\texttt{whyask\_with\_willingness.ipynb}):
\begin{verbatim}
@memo
def S1[_u: U, _g: goals](f_ctx, w_ctx, alpha):
    speaker: knows(_g)
    speaker: given(mu in mu_vals, wpp=f_prior(mu, f_ctx))
    speaker: given(t in theta_vals, wpp=1)
    speaker: given(v in v_vals, wpp=v_prior(v, w_ctx))   # <-- new
    speaker: given(phi in phi_vals, wpp=1)                # <-- new
    speaker: wants(payoff = eu_s1(u, _g, mu, t, v, phi, f_ctx))
    speaker: chooses(u in U, wpp=exp(alpha * EU[payoff]))
    return Pr[speaker.u == _u]
\end{verbatim}

The willingness version adds two \texttt{given} lines and passes $(v, \phi)$ to
the utility function. The \texttt{memo} system handles marginalization over all
four latent variables automatically.

% ============================================================
\section{Parameters}
% ============================================================

\begin{center}
\renewcommand{\arraystretch}{1.4}
\begin{tabular}{lll}
\toprule
\textbf{Parameter} & \textbf{Value} & \textbf{Role} \\
\midrule
$\alpha$           & 10   & softmax rationality \\
$\delta$           & 0.7  & face-threat cost of \textsc{imp} \\
$k$                & 10   & Beta prior concentration for $\mu$ (and $v$) \\
$c_{\text{can}}$   & 0.1 (notebook) / 0 (minimal) & question cost \\
$f_{\text{ctx}}$   & varies & ability context (key input) \\
$w_{\text{ctx}}$   & varies (willingness model only) & willingness context \\
\bottomrule
\end{tabular}
\end{center}

Same three free parameters ($\alpha, \delta, k$) in both versions.
The willingness extension adds $w_{\text{ctx}}$ as a context specification, not
a fit parameter.

% ============================================================
\section{Derivation Chain}
% ============================================================

\paragraph{Ability-only.}
\[
  f_{\text{ctx}}
  \xrightarrow{\mathrm{Beta}}
  \pi(\mu)
  \xrightarrow{\mathrm{MI}}
  \mathrm{IG}
  \xrightarrow{\mathrm{EU}}
  P_{S1}
  \xrightarrow{\mathrm{Bayes}}
  P_{A1}(g)
  \xrightarrow{\mathrm{EU}}
  P_{A1}(\textsc{act})
  \xrightarrow{\mathrm{EU}}
  P_{S2}
  \xrightarrow{\mathrm{Bayes}}
  P_{A2}(g)
  \xrightarrow{\mathrm{EU}}
  P_{A2}(\textsc{act})
\]

\paragraph{Willingness extension.} Same chain, but:
\[
  (f_{\text{ctx}}, w_{\text{ctx}})
  \xrightarrow{\mathrm{Beta}\times\mathrm{Beta}}
  \pi(\mu)\times\pi(v)
  \;\longrightarrow\; \cdots
\]
IG still depends only on $\pi(\mu)$; the willingness prior $\pi(v)$ enters only
through the comply term in the speaker and response utilities.

\bigskip
\begin{thebibliography}{9}
\bibitem{frank2012predicting}
Frank, M.\ C., \& Goodman, N.\ D.\ (2012).
Predicting pragmatic reasoning in language games.
\textit{Science}, 336(6084), 998.

\bibitem{goodman2016pragmatic}
Goodman, N.\ D., \& Frank, M.\ C.\ (2016).
Pragmatic language interpretation as probabilistic inference.
\textit{Trends in Cognitive Sciences}, 20(11), 818--829.

\end{thebibliography}

\end{document}
